% ---
% Reemplaza líneas 2666–2719 de Docs/Technical_Report.tex.
% Decisiones que lo respaldan: D-006, D-008, D-009, D-023, H-007.
% ---

\subsection{Módulo 3: Extracción de características}
\label{sec:modulo3}

El módulo de extracción transforma cada par (imagen, máscara) producido por M2 en un vector
numérico $\mathbf{x}\in\mathbb{R}^{67}$ que describe la morfología de la lesión. Se emplea
PyRadiomics \cite{vanGriethuysen2017} sobre la imagen original, sin filtros, con cuatro familias de
características:

\begin{itemize}
  \item \textit{Primer orden} (18): estadísticas de la distribución de intensidades dentro de la
        lesión, como la media, los percentiles, la energía, la asimetría y la curtosis.
  \item \textit{Forma 2D} (9): descriptores geométricos del contorno, como el perímetro, los
        ejes mayor y menor, el diámetro máximo, la elongación y la esfericidad.
  \item \textit{GLCM} (24): textura a partir de la co-ocurrencia de pares de niveles de gris.
  \item \textit{GLRLM} (16): textura a partir de la longitud de las secuencias de píxeles con el
        mismo nivel de gris.
\end{itemize}

% respaldo: «In the case that no re-segmentation range may be defined due to arbitrary intensity units (e.g. raw MRI data and many spatial filters), the use of the fixed bin size discretisation algorithm is not recommended.» (Zwanenburg2020IBSI, p. 24)
% respaldo: «whereas grey level co-occurence matrix (GLCM) features may be better modeled with a higher number of grey levels (e.g. 32 or 64)» (Zwanenburg2020IBSI, p. 26)
Las intensidades se discretizan en un número fijo de 32 niveles (\texttt{binCount}), no en un
ancho fijo. IBSI recomienda el número fijo cuando las unidades de intensidad son arbitrarias, como
en la película digitalizada del CBIS-DDSM, y desaconseja el ancho fijo cuando no puede definirse un
rango absoluto \cite{Zwanenburg2020IBSI}. El manual no fija el número de niveles; 32 está entre los
que menciona como ejemplo para la GLCM, y no se optimizó sobre los datos para no ajustar el
preprocesamiento a las etiquetas. Con un ancho fijo, los datos de 16 bits generaban más de mil niveles. Las matrices de textura usan distancia 1 y
promedian las cuatro direcciones del plano. Como las mamografías no traen espaciado de píxel, las
características de forma se expresan en píxeles.

La extracción se aplica a las 3\,568 lesiones y a la variante con CLAHE de M2, sin fallos ni
valores faltantes. Masas y calcificaciones se extraen y se guardan por separado.

Como alternativa se contempló un codificador convolucional preentrenado
(ResNet-18 o EfficientNet-B0). Se eligió PyRadiomics porque cada característica tiene un significado morfológico
definido, lo que permite interpretar qué propiedades de la lesión conserva o pierde cada
transformación.
